\documentclass[10pt,twocolumn,letterpaper]{article}
\usepackage[review]{cvpr}
\usepackage[utf8]{inputenc}
\usepackage{amsmath,amssymb,amsfonts}
\usepackage{bm}
\usepackage{booktabs}
\usepackage{graphicx}
\usepackage{microtype}
\usepackage{xcolor}
\usepackage{tabularx}
\usepackage{multirow}
\definecolor{cvprblue}{rgb}{0.21,0.49,0.74}
\usepackage[pagebackref,breaklinks,colorlinks,allcolors=cvprblue]{hyperref}
\backrefparscanfalse
\def\paperID{*****}
\def\confName{ICCV}
\def\confYear{2027}
\title{\textbf{Supplementary Material: Map Growth and Policy Routing in Budgeted Online RGB-D 3DGS}}
\author{Anonymous ICCV Submission\\Paper ID: *****}
\renewcommand{\thetable}{S\arabic{table}}
\renewcommand{\thefigure}{S\arabic{figure}}
\newcommand{\oursPSNR}{27.74}
\newcommand{\oursStd}{0.03}
\newcommand{\gainBest}{2.36}
\newcommand{\gainRTG}{4.21}
\newcommand{\gainErr}{3.53}
\newcommand{\gainNoop}{8.89}
\newcommand{\gapFull}{0.15}
\newcommand{\headNoop}{98.3}
\newcommand{\headBest}{94.0}
\newcommand{\mapRatio}{3.2}
\newcommand{\covErrA}{3.53}
\newcommand{\covErrB}{1.71}
\newcommand{\budgA}{3.69}
\newcommand{\budgB}{2.94}
\newcommand{\matOffA}{0.38}
\newcommand{\matOnA}{0.87}
\newcommand{\propMin}{1.56}
\newcommand{\propMax}{3.73}
\newcommand{\growMin}{60}
\newcommand{\growMax}{87}
\newcommand{\errPSNR}{24.21}
\newcommand{\oursMap}{22}
\newcommand{\errMap}{72}

\begin{document}
\maketitle
\appendix
\section{Short-horizon ablations}
\begin{figure}[tbp]
\centering
\includegraphics[width=\columnwidth]{figs/fig6_budget.pdf}
\caption{\textbf{Budget curve} (fr2\_xyz, 30 frames$\times$3, matched modeled budgets): policies are similar across the tested budget range at this short horizon. This experiment does not establish where the 150-frame policies separate as a function of budget.}
\label{fig:budget}
\end{figure}
\begin{figure}[tbp]
\centering
\includegraphics[width=\columnwidth]{figs/fig3_kappa_ablation.pdf}
\caption{\textbf{Replica office0 ablation} (30 frames, $n=3$). Disabling coverage throttling changes mean PSNR by $0.82\,\mathrm{dB}$; switching from temporal$\times$influence to raw-error changes it by a further $4.69\,\mathrm{dB}$. These are descriptive short-run differences without a reported significance test.}
\label{fig:kappa}
\end{figure}

\section{Historical and Cross-System Context}
The following tables use different substrates or protocols from the corrected matched-budget benchmark.
\textbf{Historical component ablation (Tab.~\ref{tab:ablation}, pose-bugged substrate, 150 frames, 5 seeds).} Relative to A0, A2 (coverage) differs by $+1.04\,\mathrm{dB}$ and A3 (coverage plus backlog) by $+1.08\,\mathrm{dB}$. The direct A3--A2 increment is $+0.04\,\mathrm{dB}$; no direct significance test is reported, and comparisons against A0 have Holm-adjusted $p=0.25$. These numbers are not transferred to the corrected benchmark.
A native local SplaTAM run on the same sequence (144f, seed 42: final $27.20$ / mean $23.03\,\mathrm{dB}$, SSIM $0.974$, LPIPS $0.046$, depth $1.26\,\mathrm{cm}$, ATE $0.41\,\mathrm{cm}$) uses a different resolution and tracking setup; these figures are contextual only and do not test compute dilution.

\begin{table}[tbp]
\centering
\footnotesize\setlength{\tabcolsep}{3pt}
\begin{tabular}{lcccc}
\toprule
\textbf{Condition} & \textbf{Mean} & \textbf{$\Delta$ vs A0} & \textbf{Map} & \textbf{Finding} \\
\midrule
A0 baseline & 16.07 & --- & 140,998 & bloated \\
A1 + warmup & 15.32 & $-0.75$ & 140,038 & harmful \\
A2 + coverage & 17.11 & $+1.04$ & 44,865 & ns vs A0 \\
A3 + backlog & 17.15 & $+1.08$ & 43,824 & $+0.04$ vs A2 \\
A4 + learned & 14.53 & $-1.54$ & 38,317 & harmful \\
\bottomrule
\end{tabular}
\caption{\textbf{Historical component ablation on the pose-bugged substrate} (5 seeds). Values and tests are not evidence for effects on the corrected benchmark. \texttt{results/step1\_ablation\_5seeds/}.}
\label{tab:ablation}
\end{table}

\begin{table*}[t]
\centering
\footnotesize\setlength{\tabcolsep}{3pt}
\begin{tabular}{lccc}
\toprule
\textbf{Method} & \textbf{Replica} & \textbf{fr2\_xyz} & \textbf{Protocol} \\
\midrule
MonoGS & 37.50 & 16.17 & full seq + refine \\
Point-SLAM & 35.17 & --- & full seq + refine \\
RTG-SLAM & 32.79 & 17.08 & full seq + refine \\
SplaTAM & --- & 24.50 & full seq + refine \\
SplaTAM (local run, 144f) & --- & 27.20 final / 23.03 mean & SLAM 640$\times$480, no refine \\
Photo-SLAM & --- & 21.07 & full seq + refine \\
Ours (error\_only) & 20--33 & 24.21 & budgeted, 100--150f \\
Ours (throttled) & 20--29 & \textbf{27.74} & budgeted, 100--150f \\
\bottomrule
\end{tabular}
\caption{\textbf{Published-SOTA context (protocol-mismatched) and one local SplaTAM run.} Published results are shown for context, not as matched comparisons. The local run uses the same sequence but a different resolution and tracking setup; details are in the text.}
\label{tab:sota}
\end{table*}

\section{Additional Evaluation}
\begin{table}[tbp]
\centering
\footnotesize\setlength{\tabcolsep}{3pt}
\begin{tabular}{lcccc}
\toprule
\textbf{Method} & \textbf{fr2\_xyz} & \textbf{fr1\_xyz} & \textbf{Horizon} & \textbf{System?} \\
\midrule
Ours (GN VO) & 0.9 & 2.0 & 100f & no (baseline) \\
MonoGS & 1.71 & --- & full & yes \\
SplaTAM & 0.50 & --- & full & yes \\
Photo-SLAM & 0.63 & --- & full & yes \\
\bottomrule
\end{tabular}
\caption{\textbf{ATE RMSE (cm).} Rows 2--4 are published full-system numbers; ours is a decoupled odometry sanity check on fixed horizons.}
\label{tab:ate}
\end{table}


\section{Further Experiments}
The following experiments would refine the mechanism and extend its scope.
\textbf{Map-growth mediation:} vary continuous spawn rates and pruning strength while holding the candidate error threshold fixed; compare trajectories at matched map size and depth quality. A budget proportional to current map size would also test whether the fourfold rescue persists under a less coarse intervention.
\textbf{Router generalization:} obtain a third sensor domain and scenes outside the reported set for threshold calibration; then test structured missing-depth artifacts and geometry-preserving TUM hole repair. The present random-hole and valid-depth-smoothing interventions did not reverse the preferred policy.
\textbf{Longer mapping evaluation:} use sequences of at least 500 frames, multiple start windows, a real tracker, and depth or geometric metrics.
\textbf{Corrected utility labels:} regenerate counterfactual labels on the corrected pose substrate before drawing conclusions about learned ranking.
\textbf{Finer-grained adaptivity:} the 7-scene routing table shows motion/exploration regime matters within a sensor (fr1\_room, fr1\_rpy trail by $\approx -1.1\,\mathrm{dB}$); extending the observable cue beyond hole-rate (e.g.\ optical-flow magnitude, coverage velocity) could route rapid-motion frames to instantaneous error while keeping temporal smoothing elsewhere.




\section{Hole-rate policy routing (auxiliary)}
Tables~S4--S6 come from earlier code revisions, before the evaluator correction (standard windowed SSIM; alpha-normalized depth). Their SSIM and depth values are not comparable with the main-paper tables, and their PSNR may differ by a few tenths of a dB.
Main-paper scene dependence motivates asking whether an observable cue can pick the better policy. We report this as an auxiliary analysis; it does not support a mechanism claim.
\textbf{Hole-rate policy routing.}
Raw depth hole-rate separates the evaluated TUM and Replica datasets: TUM $0.18$--$0.31$ vs Replica $\approx 0.000$ (300--500$\times$). Since the datasets also differ in rendering, pose quality, and scene content, this association does not establish that sensor noise causes the policy ranking to change.
The score $n = \mathrm{clip}(2.5 \cdot \mathrm{holes} + 1.0 \cdot \mathrm{edge})$ routes $n \le 0.10 \to$ raw-error selection ($\kappa = 0.98$) and $n \ge 0.30 \to$ temporal$\times$influence + throttling ($\kappa = 0.90$), computed pre-hygiene in milliseconds on CPU.
Formally, with $h_t$ the raw missing-depth rate and $e_t$ the raw edge density of frame $t$:
{\footnotesize
\begin{equation}
n_t = \min(1,\, 2.5 h_t + 1.0 e_t), \quad
\pi_t = \begin{cases}
\texttt{err\_only} & n_t \le 0.10 \\
\text{blend} & 0.10 < n_t < 0.30 \\
\texttt{temporal} & n_t \ge 0.30
\end{cases}
\end{equation}
}
with coverage knee $\kappa \in \{0.98, 0.90\}$ tied to the chosen arm.
The score and thresholds were chosen after inspecting five development scenes (Fig.~\ref{fig:separation}), including \texttt{fr1\_xyz}. That scene is therefore development evidence for the router, although its separate 150-frame mapping-policy comparison uses frozen settings. A seven-scene audit is reported separately in Tab.~\ref{tab:routing7}; scene-level independence must be tracked separately from seed-level repeats.
Two earlier features failed held-out separation and are reported as such: post-hygiene invalid ratio (conflates far-but-valid geometry --- 18\% of Replica room0 exceeds $4\,\mathrm{m}$ --- with holes) and Laplacian texture energy (measures sharpness, inverted: crisp synthetic scores \textit{higher}).
Operating on pre-hygiene raw depth was the fix (Fig.~\ref{fig:separation}).
A mixed-regime test injects 30\% holes mid-trajectory on Replica and the router changes arms on 49/50 frames at the transition. This verifies feature-triggered switching in that test; end-to-end PSNR trails fixed \texttt{error\_only} by $0.85\,\mathrm{dB}$, so a quality benefit on mixed real sequences is not established.

\begin{figure}[tbp]
\centering
\includegraphics[width=\columnwidth]{figs/fig1_noise_separation.pdf}
\caption{\textbf{Depth score on five development scenes.} The five plotted scenes include \texttt{fr1\_xyz}; its mapping-policy comparison is distinct from independent router validation. The seven-scene router audit appears in Tab.~\ref{tab:routing7}. The separation shown here measures dataset association, not sensor-noise causality.}
\label{fig:separation}
\end{figure}

\begin{table}[tbp]
\centering
\footnotesize\setlength{\tabcolsep}{3pt}
\begin{tabular}{lcccc}
\toprule
\textbf{Policy} & \textbf{Mean} & \textbf{Final} & \textbf{SSIM} & \textbf{Map} \\
\midrule
Fixed OURS & 26.85 & 26.07 & 0.9244 & 78,825 \\
Routed (err) & 30.17 & 28.02 & 0.9584 & 115,964 \\
\midrule
$\Delta$ & +3.32 $\checkmark$ & +1.95 $\checkmark$ & +0.034 & --- \\
\bottomrule
\end{tabular}
\caption{\textbf{Adaptive validation, Replica \texttt{office0}, 99 frames $\times$ 8 seeds.} Router picks \texttt{error\_only} (holes $\approx 0.02$): $p = 0.0078$, $d_z = 71.9$. No fixed policy wins both substrates.}
\label{tab:adaptive}
\end{table}

\begin{table}[tbp]
\centering
\footnotesize\setlength{\tabcolsep}{2pt}
\begin{tabular}{lcccc}
\toprule
\textbf{Scene} & \textbf{Level} & \textbf{Error} & \textbf{Temporal} & \textbf{$\Delta$} \\
\midrule
Replica office0 & 0\% & 30.18 & 26.84 & $-3.33$ \\
 & 15\% & 30.32 & 26.75 & $-3.58$ \\
 & 30\% & 30.31 & 26.62 & $-3.69$ \\
\midrule
TUM fr2\_xyz & 0 & 24.12 & 27.69 & $+3.57$ \\
 & 0.5 & 24.10 & 27.64 & $+3.54$ \\
 & 1.0 & 24.21 & 27.62 & $+3.41$ \\
\bottomrule
\end{tabular}
\caption{\textbf{Within-scene depth interventions} (150 frames on TUM, 99 on Replica, $\times$ 3 seeds). $\Delta$ is temporal$\times$influence with coverage minus unthrottled raw error, measured on identical perturbed frames. Replica also receives Gaussian depth noise with standard deviation $0.05$ times the hole fraction in meters. TUM smoothing acts only on valid depth after the standard hygiene stage; it does not remove remaining holes. No level reverses policy ordering. Full paired intervals: \texttt{results/depth\_intervention/analysis.json}.}
\label{tab:depth_intervention}
\end{table}

\begin{table}[tbp]
\centering
\footnotesize\setlength{\tabcolsep}{3pt}
\begin{tabular}{lccc}
\toprule
\textbf{Scene} & \textbf{Routed PSNR} & \textbf{Fixed reference} & \textbf{Route} \\
\midrule
fr2\_xyz 150f & 27.66 & 27.64 (ours) & temporal \\
office0 99f & 30.36 & 30.17 (err) & error\_only \\
office1 100f & 33.11 & 33.48 (err) & error\_only \\
room1 100f & 24.17 & 24.29 (err) & error\_only \\
fr3\_static 100f & 25.95 & 25.96 (ours) & temporal \\
fr1\_room 100f & 15.50 & 16.54 (err) & temporal \\
fr1\_rpy 100f & 19.93 & 21.06 (err) & temporal \\
\bottomrule
\end{tabular}
\caption{\textbf{7-scene router audit.} Routed PSNR uses $n=3$; fixed values are unpaired references. Office0 uses $n=3$ routed vs $n=8$ fixed. Five scenes are held out; two overlap development.}
\label{tab:routing7}
\end{table}

\textbf{Policy choice differs across the two evaluated substrates (Tab.~\ref{tab:adaptive}).}
On synthetic Replica the fixed combined policy loses $3.32\,\mathrm{dB}$; routing to raw-error selection recovers $3.32$ [$3.29$, $3.35$], $p = 0.0078$.
\textbf{Seven-scene routing (Tab.~\ref{tab:routing7}, exploratory $n=3$, no scene labels):} routed PSNR is close to fixed-policy reference values on five scenes and about $1.1\,\mathrm{dB}$ lower on exploration-heavy \texttt{fr1\_room} and rapid-rotation \texttt{fr1\_rpy}. These are unpaired, exploratory comparisons; the table reports references rather than inferential deltas.
Domain dependence is thus finer than real-vs-synthetic: motion and exploration regime matter within a sensor, and the current router does not capture them --- reported as-is.
\textbf{Depth intervention (Tab.~\ref{tab:depth_intervention}).} On Replica \texttt{office0}, raising random missing-depth rate to 30\% with accompanying $1.5\,\mathrm{cm}$ Gaussian noise makes the frozen score exceed its temporal-route threshold, yet raw-error selection remains ahead by $3.69$ $[3.62,3.78]$ dB across three seeds. On TUM \texttt{fr2\_xyz}, bilateral smoothing of valid depth after the standard hygiene stage changes the temporal-minus-error advantage only from $+3.57$ to $+3.41$ dB. These interventions do not support a causal claim that missing depth alone determines the better policy. The Replica masks are independent pixel noise rather than a full Kinect artifact model; the TUM intervention leaves missing pixels masked. Thresholds were not recalibrated on independent scenes because every available TUM and Replica scene appears in the reported scene set, and no third sensor domain is present locally.

\section{Counterfactual utility and negative results (auxiliary)}
\textbf{Measurement instrument.} Per-frame scheduling uses $C_i(K) = C_i^{\mathrm{base}} + K \cdot C_i^{\mathrm{step}}$ and $\Delta Q$ measured by exact full-scene counterfactual renders (active edits embedded in all $N$ splats, avoiding depth-interleaving gradient bias of affine-background approximations).
The oracle $U_i^\star = \Delta Q_i^{\mathrm{exact}} / \Delta C_i$ is our measurement instrument, never a runtime input.
It reveals: (i) negative utility in $2.7$--$38.7\%$ of interventions by stratum (pose-fixed regeneration), with only 2/75 negative interventions at depth discontinuities; (ii) interaction is protocol-dependent --- corrected joint-group measurements give $R_{\mathrm{add}}(4) = 0.95$, $R_{\mathrm{add}}(16) = 0.82$ (relatively additive), while a separate pairwise protocol shows stronger interaction; pointwise top-$K$ is therefore treated as a reference, not a bound, without a universality claim; (iii) offline TwoHeadMLP rank signal ($\rho \approx 0.20$, NDCG@20 $0.46$, measured on pre-pose-fix labels whose absolute values are polluted but whose within-set method comparison remains valid) reverses to $-2.51\,\mathrm{dB}$ closed-loop. Conditional-oracle selection numbers are withheld pending a rerun under fork-correct semantics.
Pose-fixed counterfactuals are negative in 29/75 flat, 16/75 texture, 14/75 edge and 2/75 depth-discontinuity interventions (candidate-level proportions from five seeds; no cluster-aware intervals). Uniform random selection from the same stratified pools has an expected negative-intervention rate of $20.3\%$ (95\% seed-cluster bootstrap CI $16.0$--$24.3\%$), so the stratum range alone does not establish ranking quality; the per-candidate labels needed to compare the selected top-error subset were not saved.
\textbf{Negative results.}
\textbf{Learned utility reverses closed-loop} ($-2.51\,\mathrm{dB}$ in the historical one-seed factorial on the pre-fix substrate; GradNorm beats the MLP pointwise $0.51 > 0.25$): offline labels ignore trajectory shift and inter-Gaussian coupling, so we ship no learned component. This result is separate from the corrected-pose factorial in Tab.~2 of the main paper.
\textbf{Contextual re-ranking buys no online gain}: adaptive greedy shows no significant improvement over static pointwise selection in direct multi-seed budget sweeps (Wilcoxon n.s., CIs span zero) while incurring $424\times$ selection overhead ($76.36$ vs $0.18\,\mathrm{ms}$). Supporting rank-stability/oracle-gap numbers predate the fork-semantics fix and are withheld pending regen.

\section{Additional main-benchmark figures}
\begin{figure}[tbp]
\centering
\includegraphics[width=\columnwidth]{figs/fig2_main_deltas.pdf}
\caption{\textbf{Paired mean-PSNR differences} on fr2\_xyz (150 frames, $n=8$). Points show OURS $-$ baseline; bars are 95\% paired-bootstrap CIs. Positive favors OURS. FULL is $\gapFull\,\mathrm{dB}$ higher; all six contrasts pass Holm correction ($p=0.0469$).}
\label{fig:deltas}
\end{figure}

\begin{figure}[tbp]
\centering
\includegraphics[width=\columnwidth]{figs/fig4_pareto.pdf}
\caption{\textbf{Quality--map-size trade-off} on fr2\_xyz (150 frames, $n=8$). Among policies that optimize, only OURS and FULL lie on the frontier; FULL gains $\gapFull\,\mathrm{dB}$ with a $\mapRatio\times$ larger map. The four growing baselines hold $\growMin$--$\growMax\mathrm{k}$ Gaussians; no-op stays compact at lower quality.}
\label{fig:pareto}
\end{figure}

\begin{figure*}[t]
\centering
\includegraphics[width=\textwidth]{figs/fig5_qualitative.pdf}
\caption{\textbf{Qualitative comparison across both evaluated substrates.} Top: an intermediate TUM \texttt{fr2\_xyz} view (90), where OURS has lower per-view error than \texttt{error\_only}; bottom: Replica \texttt{office0} view 98, where the ranking reverses. Columns show ground truth, OURS, \texttt{error\_only}, FULL, and signed RGB $L_1$ error gap (red: OURS has lower error). Insets mark the same crop in each row. This figure comes from an earlier code revision. PSNR labels are per-view values from seed 42, shown only to illustrate the visual comparison; they are not sequence means.}
\label{fig:qual}
\end{figure*}

\section{Per-scene breakdown of the combined policy (earlier code revision)}
\begin{table*}[t]
\centering
\small
\begin{tabular}{lcccccc}
\toprule
\textbf{Scene} & \textbf{Frames} & \textbf{OURS} & \textbf{\texttt{error\_only}} & \textbf{$\Delta$} & \textbf{Map (k)} & \textbf{$n$} \\
\midrule
fr2\_xyz (TUM) & 150 & \oursPSNR & \errPSNR & $+\gainErr$ $\checkmark$ & \oursMap / \errMap & 8 \\
fr1\_xyz (TUM) & 150 & 23.47 & 21.09 & $+2.38$ $\checkmark$ & 39 / 81 & 8 \\
fr1\_desk (TUM) & 150 & 22.55 & 22.34 & $+0.21$ & 27 / 56 & 3 \\
fr1\_room (TUM) & 100 & 15.53 & 16.54 & $-1.01$ & 92 / 97 & 3 \\
fr1\_rpy (TUM) & 100 & 19.92 & 21.06 & $-1.14$ & 45 / 62 & 3 \\
fr3\_static (TUM) & 100 & 25.96 & 25.88 & $+0.08$ & 13 / 18 & 3 \\
office0 (Rep.) & 99 & 26.85 & 30.17 & $-3.32$ $\checkmark$ & 79 / 116 & 8 \\
office1 (Rep.) & 100 & 29.25 & 33.46 & $-4.21$ $\checkmark$ & 42 / 74 & 8 \\
office2 (Rep.) & 99 & 22.08 & 24.67 & $-2.59$ & 70 / 117 & 3 \\
office3 (Rep.) & 99 & 22.45 & 24.47 & $-2.02$ & 77 / 150 & 3 \\
office4 (Rep.) & 99 & 23.58 & 25.93 & $-2.36$ & 81 / 141 & 3 \\
room0 (Rep.) & 99 & 20.85 & 22.78 & $-1.93$ & 82 / 150 & 3 \\
room1 (Rep.) & 100 & 20.42 & 24.29 & $-3.87$ & 47 / 59 & 3 \\
room2 (Rep.) & 99 & 22.86 & 25.88 & $-3.03$ & 78 / 146 & 3 \\
\bottomrule
\end{tabular}
\caption{\textbf{Per-scene combined-policy breakdown, 14 sequences.} $\checkmark$ marks the reported significant comparisons at $n=8$; other entries are exploratory. The combined policy wins the two confirmatory TUM comparisons but loses on both confirmatory Replica comparisons; remaining scenes have exploratory sample sizes.}
\label{tab:scenes}
\end{table*}

\end{document}
